\documentclass[10pt,landscape]{article}
\usepackage{amssymb,amsmath,amsthm,amsfonts}
\usepackage{multicol,multirow}
\usepackage{calc}
\usepackage{ifthen}
\usepackage[landscape]{geometry}
\usepackage[colorlinks=true,citecolor=blue,linkcolor=blue]{hyperref}
\usepackage{datetime2}
\DTMsetdatestyle{iso}

\geometry{a4paper,
    top=0.5cm,
    left=0.5cm,
    right=0.5cm,
    bottom=0.5cm
    }

\pagestyle{empty}

\makeatletter
\renewcommand{\section}{\@startsection{section}{1}{0mm}%
                                {-1ex plus -.5ex minus -.2ex}%
                                {0.5ex plus .2ex}%x
                                {\normalfont\normalsize\bfseries}}
\renewcommand{\subsection}{\@startsection{subsection}{2}{0mm}%
                                {-1ex plus -.5ex minus -.2ex}%
                                {0.5ex plus .2ex}%
                                {\normalfont\small\bfseries}}
\renewcommand{\subsubsection}{\@startsection{subsubsection}{3}{0mm}%
                                {-1ex plus -.5ex minus -.2ex}%
                                {0.5ex plus .2ex}%
                                {\normalfont\small\bfseries}}
\makeatother
\setcounter{secnumdepth}{0}
\setlength{\parindent}{0pt}
\setlength{\parskip}{0pt plus 0.5ex}
% -----------------------------------------------------------------------

\title{MAST90104 Part 2 Cheat Sheet for Enoch Ko (147938)}

\begin{document}

\raggedright
\footnotesize
%Jack's Cheatsheet%
\begin{multicols}{2}
\setlength{\premulticols}{1pt}
\setlength{\postmulticols}{1pt}
\setlength{\multicolsep}{1pt}
\setlength{\columnsep}{2pt}
\section{General Linear Hypothesis (GLH) Test — Quick Sheet}

\textbf{Model \& Hypothesis.}
\[
y=X\beta+\varepsilon,\quad \varepsilon\sim(0,\sigma^2 I_n),\qquad
H_0:\;L\beta=\delta,\; L\in\mathbb R^{r\times p},\; \mathrm{rank}(L)=r
\]

\textbf{(1) Put the constraint in estimates.}
\[
\text{Gap vector: } g:=L\hat\beta-\delta
\]

\textbf{(2) Reparameterise if $X$ is not full rank.} Pick full-rank $C$ so that $\tilde X:=XC$ has independent columns and $X\beta=\tilde X\gamma$. Rewrite $H_0$ as $\tilde L\gamma=\delta$ and proceed with tildes.

\textbf{(3) Gram matrix.}
\[
G:=X^\top X\quad (\text{or } \tilde G:=\tilde X^\top \tilde X)
\]

\textbf{(4) Hat matrix \& residual sum of squares.}
\[
H:=X(X^\top X)^{-1}X^\top,\quad \hat y=Hy,\quad
SS_{\mathrm{Res}}:=y^\top(I-H)y,\quad \mathrm{MSE}:=\frac{SS_{\mathrm{Res}}}{n-p}
\]

\textbf{(5) Constraint on $\hat\beta$.}
\[
\hat\beta=(X^\top X)^{-1}X^\top y,\qquad
S:=L(X^\top X)^{-1}L^\top
\]

\textbf{(6) ``Write half the top first'' then assemble the quadratic form.}
\[
\text{Top (numerator part)}:\quad Q:=g^\top S^{-1}g
\]

\textbf{(7) Test statistic \& decision.}
\[
F=\frac{(Q/r)}{\mathrm{MSE}}=\frac{g^\top S^{-1}g/r}{SS_{\mathrm{Res}}/(n-p)}
\ \ \sim\ F_{r,\,n-p}\ \text{ under }H_0;
\quad \text{reject if }F\ge F_{1-\alpha;\,r,\,n-p}.
\]

\textbf{Two-RSS (equivalent) form.} Fit restricted ($H_0$) and full models:
\[
F=\frac{\big(SS_{\mathrm{Res}}^{H_0}-SS_{\mathrm{Res}}^{H_1}\big)/r}{SS_{\mathrm{Res}}^{H_1}/(n-p)}
\ \ \sim\ F_{r,\,n-p}.
\]

\textbf{Notes.} $X^\top X$ is the Gram matrix; $\mathrm{Var}(\hat\beta)=\sigma^2(X^\top X)^{-1}$. Reparameterisations that make $\tilde X$ full rank leave $\hat y$, $SS_{\mathrm{Res}}$, MSE, and the $F$ test unchanged.

\medskip
\section{Coding Choices for One-Way ANOVA (High-School Intuition)}
We start from the over-parameterised model with intercept:
\[
\beta=\begin{bmatrix}\mu\\ \tau_1\\ \vdots\\ \tau_k\end{bmatrix},\quad
\text{and reparameterise } \beta=C\,\gamma \text{ so } \tilde X=XC \text{ is full rank.}
\]
Use the coding that matches how you want to \emph{explain} results.

\subsection{Treatment (Reference) Coding: ``each group vs a baseline''}
Pick a baseline level (say level $k$). Keep intercept and $k-1$ group effects vs baseline.
\[
\gamma=\begin{bmatrix}\mu\\ \theta_1\\ \vdots\\ \theta_{k-1}\end{bmatrix},\qquad
C_{\text{treat}}=
\begin{bmatrix}
1 & 0      & \cdots & 0 \\
0 & 1      & \cdots & 0 \\
\vdots & \vdots & \ddots & \vdots \\
0 & 0      & \cdots & 1 \\
0 & 0      & \cdots & 0
\end{bmatrix}_{(k+1)\times k},
\quad
\beta=
\begin{bmatrix}\mu\\ \theta_1\\ \vdots\\ \theta_{k-1}\\ 0\end{bmatrix}.
\]
\textbf{Concrete for $k=3$:}
\[
C_{\text{treat}}=
\begin{bmatrix}
1&0&0\\
0&1&0\\
0&0&1\\
0&0&0
\end{bmatrix},\quad
\gamma=(\mu,\theta_A,\theta_B),\quad
\theta_A=\tau_A,\ \theta_B=\tau_B,\ \tau_C=0.
\]
\textit{How you'd say it (scores $A=70,B=80,C=90$ with $C$ baseline):}
Intercept $\approx 90$; $A$ is $-20$ vs $C$, $B$ is $-10$ vs $C$.

\subsection{Sum (Effect) Coding: ``each group vs the overall average''}
Enforce $\sum_{i=1}^k\tau_i=0$. Keep $k-1$ independent effects; the last is implied.
\[
\gamma=\begin{bmatrix}\mu\\ \alpha_1\\ \vdots\\ \alpha_{k-1}\end{bmatrix},\qquad
C_{\text{sum}}=
\begin{bmatrix}
1 & 0      & \cdots & 0 \\
0 & 1      & \cdots & 0 \\
\vdots & \vdots & \ddots & \vdots \\
0 & 0      & \cdots & 1 \\
0 & -1     & \cdots & -1
\end{bmatrix}_{(k+1)\times k},
\quad
\beta=
\begin{bmatrix}
\mu\\ \alpha_1\\ \vdots\\ \alpha_{k-1}\\ -\sum_{j=1}^{k-1}\alpha_j
\end{bmatrix}.
\]
\textbf{Concrete for $k=3$:}
\[
C_{\text{sum}}=
\begin{bmatrix}
1&0&0\\
0&1&0\\
0&0&1\\
0&-1&-1
\end{bmatrix},\quad
\gamma=(\mu,\alpha_A,\alpha_B),\ 
\tau_C=-(\alpha_A+\alpha_B).
\]
\textit{How you'd say it (scores $70,80,90$):}
Grand mean $=80$ (intercept). $A$ is $-10$ from average, $B$ is $0$, so $C$ is $+10$ (they must sum to $0$).

\subsection{Basic (Cell-Means) Coding: ``parameters are the group means''}
Let parameters \emph{be} the group means $m_i=\mu+\tau_i$; no baseline, no sums.
\[
\gamma=\begin{bmatrix}m_1\\ \vdots\\ m_k\end{bmatrix},\qquad
\bar m=\frac{1}{k}\sum_{i=1}^k m_i,\quad
\mu=\bar m,\ \tau_i=m_i-\bar m,
\]
\[
C_{\text{basic}}=
\begin{bmatrix}
\frac{1}{k}\mathbf 1_k^\top\\[3pt]
I_k-\frac{1}{k}\mathbf 1_k\mathbf 1_k^\top
\end{bmatrix}\in\mathbb R^{(k+1)\times k},
\quad
\beta=
\begin{bmatrix}
\bar m\\ m_1-\bar m\\ \vdots\\ m_k-\bar m
\end{bmatrix}.
\]
\textbf{Concrete for $k=3$:}
\[
C_{\text{basic}}=
\begin{bmatrix}
\frac{1}{3}&\frac{1}{3}&\frac{1}{3}\\
\frac{2}{3}&-\frac{1}{3}&-\frac{1}{3}\\
-\frac{1}{3}&\frac{2}{3}&-\frac{1}{3}\\
-\frac{1}{3}&-\frac{1}{3}&\frac{2}{3}
\end{bmatrix},\quad \gamma=(m_A,m_B,m_C).
\]
\textit{How you'd say it (scores $70,80,90$):}
The parameters are simply the group means: $A=70,\ B=80,\ C=90$.
\end{multicols}
\end{document}